\section{Convergence of Stochastic Processes}
\label{sec:convergence-process}

In~\Cref{sec:dynamics}, we deal
with convergence in distribution of
stochastic processes indexed by countably
infinite set, intended as weak convergence in the
product topology. Equivalently, this
means that every finite-dimensional
marginal converges in distribution.

\begin{definition}
    Let $\calA$ be a countable
    set. For random variables $(\bx^{(n)})_{n\ge 1}$ and
    $\bx^\infty$ taking values in $\R^\calA$, we say
    that $\bx^{(n)}$ converges in distribution to
    $\bx^\infty$ and write
    \[
        \bx^{(n)}\tod \bx^\infty
    \]
    if, for every $k\ge 1$ and $\alpha_1, \ldots,\alpha_k\in \calA$, we have
    \[
        (x^{(n)}_{\alpha_1},\ldots,x^{(n)}_{\alpha_k})\tod (X^\infty_{\alpha_1},\ldots X^{\infty}_{\alpha_k})\,.
    \]
\end{definition}

\noindent To show convergence in distribution, we will use the method of moments~\cite[Theorems 29.4, 30.1, 30.2]{billingsleyProbabilityBook}.
The following theorem follows from
Carleman's conditions on 
moment-determinacy of a distribution on
$\R$, combined with~\cite{petersenEquivalence}.

\begin{theorem}[Method of moments]\label{cor:bounded-limit-existence}
    Let $(\bx^{(n)})_{n\ge 1}$ be a sequence
    of stochastic processes indexed by a
    countable set $\calA$. Assume
    that 
    \begin{enumerate}
        \item All joint moments converge: for any
        $k\ge 1$ and $\alpha_1, \ldots, \alpha_k\in \calA$,
        the limit of the joint moments
    \begin{align}
        \lim_{n\to\infty} \E \left[\prod_{i=1}^k x^{(n)}_{\alpha_i}\right]\label{eq:joint-moments}
    \end{align} 
    exists.
        \item All marginals are subexponential: for every $\alpha\in\calA$, there exists $C_\alpha>0$ such that
        for all $p\ge 1$,
        \begin{align}
            \lim_{n\to\infty} \E \left(x^{(n)}_\alpha\right)^{2p}\le (C_\alpha p)^{2p}\,.\label{eq:subexp-growth-condition}
        \end{align}
    \end{enumerate}
    Then $\bx^{(n)}$ converges in distribution to the unique law on $\R^{\cal A}$ with
    moments given by~\cref{eq:joint-moments}.
\end{theorem}



\begin{lemma}[Truncation]\label{lem:truncation}
    Let $(x_n)_{n\ge 1}$ and $(y_n)_{n\ge 1}$
    be sequences of random variables such that
    \begin{enumerate}
        \item For any $K>0$, conditionally
        on $|x_n|\le K$, $(y_n)_{n\ge 1}$ converges
        in distribution.
        \item $(x_n)_{n\ge 1}$ is tight, i.e., $\sup_{n\ge 1} \Pr(|x_n|>K)\underset{K\to\infty} \longrightarrow 0$.
    \end{enumerate}
    Then, $(y_n)_{n\ge 1}$ converges in distribution.
\end{lemma}

\begin{proof}
    First, we prove:

    \begin{claim}\label{claim:y-tight}
        $(y_n)_{n\ge 1}$ is tight.
    \end{claim} 

    \begin{proof}
        For any $K,L>0$, we have $\Pr(|y_n|>L)\le \Pr(|y_n|>L \mid |x_n|\le K) + \Pr(|x_n|> K)$. Pick $K$ large enough so that the second term is bounded by $\eps$ uniformly in $n$. $(y_n)_{n\ge 1}$ is tight conditionally on $|x_n|\le K$, so there exists $L>0$ large enough so 
        that the first term is also bounded by $\varepsilon$ uniformly in $n$.
    \end{proof}

    \noindent By~\Cref{claim:y-tight} and Prokhorov's
    theorem, it remains to show that
    every subsequence of $(y_n)_{n\ge 1}$ that converges in distribution, converges
    to the same limit.
    Fix $f:\R\to \R$ to be a bounded continuous
    function and $\eps>0$. Then, by the law of total
    expectations, for any $n\ge 1$,
    \begin{align*}
        \left|\E f(y_n) - \E \left[f(y_n) \mid |x_n|\le K\right]\right| &=  \Pr(|x_n|> K)\left(\E \left[f(y_n) \mid |x_n|> K\right] - \E \left[f(y_n) \mid |x_n|\le K\right]\right)\\
        &\le 2 \|f\|_\infty \Pr(|x_n|> K)\\
        &\le \varepsilon
    \end{align*}
    by setting $K = K(\eps)$ to be a large enough constant (with the second assumption). By the first assumption, there exists $N\ge 1$ such that
    for any $n,m\ge N$,
    \[
        \left|\E \left[f(y_n) \mid |x_n|\le K\right] - \E \left[f(y_m) \mid |x_m|\le K\right]\right|\le \eps\,.
    \]
    In turn, this implies $\left|\E f(y_n) - \E f(y_m)\right|\le 3\eps$ by the triangle inequality, so $(\E f(y_n))_{n\ge 1}$
    is a Cauchy sequence, so it converges
    as $n\to\infty$. This implies
    that every weak subsequential limit of
    $(y_n)_{n\ge 1}$ converges to the same limit,
    which concludes the proof.
\end{proof}


\subsection{Connection with convergence of the empirical distribution}

Let us also remark on certain details concerning modes of convergence that are important to the use and interpretation of~\Cref{thm:convergence-distribution}.

Recall that we ``stack'' the $\bz_{\alpha}(\bA)$ for $\alpha\in\calA_1$ into a single vector with more complicated entries, $\bz_{\cA_1}(\bA) \in (\R^{\cA_1})^n$.
Using our notation from~\Cref{sec:intro}, we then sample a random coordinate of this vector, forming a further random countably infinite vector $\samp(\bz_{\cA_1}(\bA)) \in \R^{\cA_1}$.
This contains the $i$th entry of each $\bz_{\alpha}(\bA)$, for a single shared randomly chosen $i \sim \Unif([n])$.
Define the infinite random vector $Z_{\cA_1}^{\infty}$ similarly. \Cref{thm:convergence-distribution} states that:
\begin{equation}
\samp(\bz_{\cA_1}(\bA^{(n)})) \xrightarrow[n \to \infty]{\text{(d)}} Z_{\cA_1}^{\infty}\,,
\label{eq:samp-conv-d}
\end{equation}
By the Cram\'{e}r-Wold theorem, this is equivalent to: for any bounded continuous function $\varphi$ and any finitely supported vector of coefficients $c_{\alpha}$,
\[ \lim_{n \to \infty} \E_{\bA} \frac{1}{n}\sum_{i = 1}^n \varphi\left(\sum_{\alpha \in \cA_1} c_{\alpha} \bz_{\alpha}(\bA)[i]\right) = \E_{\bA} \varphi\left(\sum_{\alpha \in \cA_1} c_{\alpha} Z_{\alpha}^{\infty} \right). \]

Alternatively, we may also make sense of this statement in terms of empirical distributions, which are just the laws of the random variables $\samp(\bx)$ discussed above.
\begin{definition}[Empirical distribution]
    For $\bm x \in \R^n$, we write $\ed(\bm x) \colonequals \frac{1}{n} \sum_{i = 1}^n \delta_{\bmx[i]}$ for the \emph{empirical distribution} of the entries of $\bm x$.
\end{definition}
\noindent
Then, $\ed(\bz_{\cA_1}(\bA))$ is a random probability measure on the space $\R^{\cA_1}$, and the random variable $\samp(\bz_{\cA_1}(\bA^{(n)}))$ is a single draw from this random probability measure.
Its law is a \emph{deterministic} probability measure on the space $\R^{\cA_1}$, which is the expectation of the random measure $\ed(\bz_{\cA_1}(\bA))$ (if $\mu$ is a random measure, then its expectation takes values $(\EE \mu)(A) = \EE[\mu(A)]$).
Thus, the above~\cref{eq:samp-conv-d} is further equivalent to the weak convergence of probability measures
\begin{equation*}
\EE \ed(\bz_{\cA_1}(\bA^{(n)})) \xrightarrow[n \to \infty]{\text{(w)}} \Law(Z_{\cA_1}^{\infty})\,.
\end{equation*}
Again by the Cram\'{e}r-Wold theorem, this is equivalent to, for any finitely supported coefficient vector of $c_{\alpha}$, having
\[ \EE \ed\left(\sum_{\alpha \in \cA_1} c_{\alpha} \bz_{\alpha}(\bA^{(n)})\right) \xrightarrow[n \to \infty]{\text{(w)}} \Law\left(\sum_{\alpha \in \cA_1} c_{\alpha} Z_{\alpha}^{\infty}\right). \]
In particular, since the output $\bx_t$ of a GFOM can be viewed in the above way, we see that the empirical distributions of $\bx_t = \bx_t(\bA)$ are related to the asymptotic states $X_t^{\infty}$ by
\[ \EE \ed(\bx_t(\bA^{(n)})) \xrightarrow[n \to \infty]{\text{(w)}} \Law(X_t^{\infty})\,. \]

Thus our results, interpreted in terms of convergence of the random empirical distributions of GFOM iterates, give convergence of the expectations of random measures.
Often it is desirable to prove stronger modes of convergence in such situations, by proving that not only do we have
\[ \lim_{n \to \infty} \E_{\bA} \frac{1}{n} \sum_{i = 1}^n \varphi(\bx_t(\bA^{(n)})[i]) = \EE \varphi(X_t^{\infty})\,, \]
but also that the random variable inside the expectation \emph{concentrates} over the randomness in $\bA$.
We do not pursue this here, because it would require introducing additional assumptions on the matrices $\bA$ involved, which may vary from application to application.
As the example discussed in~\Cref{rem:cdm-factorization} shows, this kind of concentration does not follow automatically from the convergence in expectation that we show.
An instructive example is the argument in \cite{bayati2015universality}, which uses similar proof techniques to ours, but, to show that the above kind of convergence also happens in $L^2$ uses a trick involving the entrywise independence of the Wigner matrices they work with (see their Proposition~5).

In our much more general setting, it seems reasonable to ask instead for the convergence in the definition of the traffic distribution in~\cref{eq:traffic-conv} to happen in a stronger mode such as $L^2$.
We leave the exploration of such conditions and the determination of which random matrix distributions they hold for to future work.
